\section{Comparative Analysis of Related Research}

\subsection{Comparative Table of Related Research Studies}

Table~\ref{tab:lit_matrix} synthesizes the core literature evaluated
during the LuxeWay design phase, mapping each work's contribution,
limitations, and specific relevance to architectural and business
decisions embedded in the platform.

% table* spans both columns in IEEEtran two-column layout
\begin{table*}[ht]
\centering
\caption{Comparative Matrix of Related Work}
\label{tab:lit_matrix}
\begin{tabular}{p{2.2cm} p{3.8cm} p{3cm} p{3cm} p{3.5cm}}
\toprule
\textbf{Research Paper} &
\textbf{Main Contribution} &
\textbf{Strengths} &
\textbf{Limitations} &
\textbf{Relation to LuxeWay} \\
\midrule
Gupta et al.\ \cite{gupta2025double} &
Modern vehicle rental platform using MERN stack. &
Structured UI/UX workflows and admin dashboards. &
Centralized MongoDB creates single points of failure. &
Foundational reference for RESTful API and role-based UI design. \\
\midrule
Madhusudan \cite{madhusudan2019applying} &
Decentralized booking/payment on Ethereum with ECIES encryption. &
High-integrity transactions and automated penalties. &
Prohibitive gas fees and confirmation latency. &
Informs escrow-based deposit logic and digital contract design. \\
\midrule
El~Hog et al.\ \cite{elhog2025decentralized} &
Hybrid CarChain architecture (BSC + Firebase). &
Optimized throughput and reduced intermediary costs. &
Residual dependency on centralized storage. &
Paradigm for minimizing transaction overhead. \\
\midrule
Soppert et al.\ \cite{soppert2024benefit} &
Economic optimization model unifying CR and CS fleet usage. &
Quantified 23\%--120\% revenue increase from hybrid models. &
Scoped to B2C fleets; P2P legal dimensions unaddressed. &
Validates the CR+CS hybrid listing strategy in LuxeWay. \\
\midrule
Neifer et al.\ \cite{neifer2024trust} &
OBD telematics-driven reputation scoring for P2P rentals. &
Objective, manipulation-resistant trust signals. &
Requires physical OBD hardware per vehicle. &
Conceptual basis for LuxeWay's multi-dimensional review model. \\
\midrule
Kauffman \& Naldi \cite{kauffman2020research} &
Regulatory and economic landscape of sharing mobility. &
Systematic mapping of sustainable growth boundaries. &
Theoretical; no concrete system design. &
Informs platform compliance posture and market positioning. \\
\bottomrule
\end{tabular}
\end{table*}

\subsection{Synthesis and Design Implications}

Three cross-cutting observations emerge from this analysis and directly
shape LuxeWay's design:

\textbf{Technology.} The contrast between blockchain-native
solutions~\cite{madhusudan2019applying, elhog2025decentralized} and
conventional web stacks~\cite{gupta2025double} reveals a trade-off
between decentralization integrity and deployment pragmatism. LuxeWay
adopts a conventional full-stack architecture (React~18 + Spring
Boot~3 + SQL Server) while incorporating the \textit{behavioral
principles} of blockchain design---immutable records, transparent audit
trails, automated contracts---without on-chain execution overhead.

\textbf{Business model.} Soppert et al.'s~\cite{soppert2024benefit}
finding that hybrid CR+CS fleet management yields 23--120\% revenue
improvement justifies LuxeWay's dual-ecosystem architecture supporting
both daily car rental and shorter-duration motorbike sharing under one
owner dashboard.

\textbf{Trust infrastructure.} The spectrum from subjective reviews to
telematics-driven scoring~\cite{neifer2024trust} identifies a practical
middle ground: a structured multi-dimensional review protocol that
increases signal richness without requiring per-vehicle hardware,
preserving accessibility for individual private owners.
